\begin{question}{25}{
    Een militaire eenheid bewaakt een strategisch logistiek knooppunt in een conflictgebied. 
    De basis wordt regelmatig doelwit van vijandelijke \emph{loitering munitions} (kamikazedrones).

    Uit operationele data over de afgelopen periode blijkt dat er gemiddeld $\lambda = 16$ van dit soort drone-aanvallen per dag plaatsvinden volgens een Poissonproces.
    
    De basis is uitgerust met een geautomatiseerd Counter-UAS (Unmanned Aircraft Systems) verdedigingssysteem. 
    Dit systeem kan per dag maximaal $20$ drones effectief uitschakelen (beperkt door munitievoorraad en herlaadtijd). 
    Aanvallen boven deze capaciteit kunnen niet meer tijdig worden geneutraliseerd en leiden tot inslagen op de infrastructuur.
}
    \subquestion{5}{
        Bereken de kans dat de basis in een willekeurige week ($7$ dagen) \emph{minimaal} $125$ drone-aanvallen registreert.
    }
    \answer{
        Laat $X$ het aantal aanvallen in een week ($7$ dagen) zijn. 
        Aangezien de meeteenheid in dagen is, geldt dat we meten over een periode van $t = 7$ meeteenheden. \rubric{1}
        Aangezien het daggemiddelde $\lambda = 16$ is, geldt voor de weekintensiteit:
        \[
            \mu = \lambda \cdot t = 16 \cdot 7 = 112. \rubric{1}
        \]
        Het aantal aanvallen per week $X$ is dus Poisson verdeeld met gemiddelde $\mu = 112$.\rubric{1}
        We zoeken de kans $P(X \geq 125)$, die we als volgt berekenen:
        \begin{align*}
            P(X \geq 125)   &= 1 - P(X \leq 124) \\
                            &= 1 - \poissoncdf(\mu=112, k=124) \\
                            &\approx 0.1199. \rubric{2}
        \end{align*}
        De kans dat er in een willekeurige week minimaal $125$ aanvallen plaatsvinden is \textbf{0.1199} (of \SI{11.99}{\percent}).
    }

    \subquestion{4}{
        Er heeft zojuist een aanval met een kamikazedrone plaatsgevonden.
        Bereken de kans dat de volgende aanval binnen twee uur plaatsvindt.
    }
    \answer{
        Laat $Y$ het aantal aanvallen zijn in de volgende twee uur ($\frac{1}{12}$ dag),
        oftewel $Y$ is Poisson verdeeld met gemiddelde $\mu = \lambda \cdot t = 16 \cdot \frac{1}{12} \approx 1.3333$.\rubric{1}
        
        De kans dat de volgende aanval binnen twee uur plaatsvindt is gelijk aan
        de kans dat er de komende twee uur minstens één aanval plaatsvindt: \rubric{1}

        \begin{align*}
            P(Y \geq 1) &= 1 - P(Y = 0) \\
                        &= 1 - \poissonpdf(\mu=1.3333, k=0) \\
                        &\approx 0.7364.\rubric{2}
        \end{align*}
    }

    \subquestion{7}{
        Bereken de kans dat er in een willekeurige maand ($30$ dagen) op minstens vier dagen sprake is van overbelasting van het verdedigingssysteem.
    }
    \answer{
        De kans dat het verdedigingssysteem op een willekeurige dag overbelast raakt is 
        \begin{align*}
            P(X > 20)   &= 1 - P(X \leq 20) \\
                        &= 1 - \poissonpdf(\mu=16, k=20) \\
                        &= 0.1318. \rubric{2}
        \end{align*}

        Laat $Z$ het aantal dagen in een maand zijn waarop het systeem overbelast raakt.
        De drone-aanvallen zijn onafhankelijk van dag tot dag, omdat ze volgens een Poissonproces plaatsvinden. \rubric{1}

        Door deze onafhankelijkheid kunnen we $Z$ modelleren als een binomiale kansvariabele met parameters $n=30$ en $p=0.1318$.\rubric{2}
            
        We zoeken de kans $P(Z \geq 4)$, die we als volgt berekenen:
        \begin{align*}
            P(Z \geq 4)     &= 1 - P(Z \leq 3) \\
                            &= 1 - \binomcdf(n=30, p=0.1318, k=3) \\
                            &\approx 0.5711. \rubric{2}
        \end{align*}
        De kans dat er in een willekeurige maand minstens vier dagen sprake is van overbelasting, is dus ongeveer \textbf{0.5711} (of \SI{57.11}{\percent}).
    }

    \subquestion{7}{
        De eenheid wil de capaciteit van het verdedigingssysteem verhogen zodat de kans op overbelasting per dag kleiner dan \SI{0.1}{\percent} ($0.001$) wordt. 
        Bereken de minimale capaciteit (het aantal vijandelijke drones dat per dag uitgeschakeld kan worden) om aan deze eis te kunnen voldoen.
    }
    \answer{
        Laat $C$ de gezochte capaciteit per dag zijn.
        Er moet gelden dat de kans op overbelasting kleiner is dan $0.001$:
        \[
            P(Y > C) = 1 - \poissoncdf(\mu=16, k=C) < 0.001 \rubric{1}
        \]
        Deze ongelijkheid kunnen we oplossen met behulp van de \emph{numerieke solver}:
        \begin{align*}
            y_1 &= 1 - \poissoncdf(\mu=16, k=X) \\
            y_2 &= 0.001 \rubric{2}
        \end{align*}
        Aangezien de capaciteit $C$ een geheel getal moet zijn, zoeken we de kleinste waarde van $C$ waarvoor geldt dat $y_1 < y_2$.\rubric{1}
        Dit doen we door een tabel te maken voor opeenvolgende waarden van $C$:
        \begin{itemize}
            \item Voor $C = 29 \implies P(Y > 29) \approx 0.0011 > 0.001$ (voldoet nog niet)\rubric{1}
            \item Voor $C = 30 \implies P(Y > 30) \approx 0.0006 < 0.001$ (voldoet wel)\rubric{1}
        \end{itemize}
        De minimale capaciteit die nodig is om aan de veiligheidseis te voldoen is gelijk aan \textbf{30 drones per dag}.\rubric{1}
    }
    
    \subquestion{2}{
        Door escalatie in het conflictgebied wordt de basis nu ook bestookt door een tweede vijandelijke eenheid, die gemiddeld $6$ drone-aanvallen uitvoert per dag volgens een Poissonproces.
        Leg uit met welke kansverdeling je nu het totale aantal aanvallen per dag kunt modelleren.
    }
    \answer{
        Het totale aantal aanvallen per dag kan worden gemodelleerd als een enkele Poissonverdeling.\rubric{1}
        De som van twee onafhankelijke Poisson-processen is namelijk zelf ook weer een Poisson-proces, waarvan de parameter gelijk is aan de som van de individuele gemiddelden ($\lambda = 16 + 6 = 22$).\rubric{1}
    }

    
\end{question}